\begin{textsample}{POS Dim 1 – sc – Score 83.00 – sc/sc\_em\_67.txt}  \label{ex:f1_pos_003}
5.2 FOTO(CIDADE)GRAFIAS EM MOVIMENTO Tema : Cidades, narrativas, memórias, fotografias e patrimônios.
Nome da \textbf{trilha} de \textbf{aprofundamento} : Foto(cidade)grafias em movimento.
Área(s) do(s) conhecimento(s) : \textbf{Ciências} Humanas e Sociais Aplicadas, Linguagens e suas \textbf{tecnologias}, Matemática e suas \textbf{tecnologias} e Ciências da Natureza e suas \textbf{tecnologias}.
Carga horária : 160 h ou 240 h, a \textbf{depender} da \textbf{matriz} curricular em funcionamento na Unidade Escolar. \textbf{Aulas} semanais : 10 ou 15 \textbf{aulas}, conforme \textbf{matriz} em funcionamento na Unidade Escolar.
Perfil docente : O corpo docente será composto por professores habilitados, com formação em licenciaturas específicas em suas áreas.
Para compor essa \textbf{trilha}, serão necessários os seguintes profissionais : - Área das Ciências Humanas e Sociais Aplicadas - Profissional do componente História ou Geografia : 2 \textbf{aulas} ou 3 \textbf{aulas}, a \textbf{depender} da \textbf{matriz} em funcionamento na Unidade Escolar. - Profissional do componente Sociologia ou Filosofia : 1 \textbf{aula} ou 2 \textbf{aulas}, a \textbf{depender} da \textbf{matriz} em funcionamento na Unidade Escolar. - Área das Linguagens e suas \textbf{tecnologias} - Profissional do componente Língua Portuguesa ou \textbf{Inglesa} : 2 \textbf{aulas} ou 3 \textbf{aulas}, a \textbf{depender} da \textbf{matriz} em funcionamento na Unidade Escolar. - Profissional do componente Arte : 2 \textbf{aulas} ou 3 \textbf{aulas}, a \textbf{depender} da \textbf{matriz} em funcionamento na Unidade Escolar. - Profissional do componente Educação Física : 1 \textbf{aula} em todas as \textbf{matrizes}. - Área da Matemática e suas \textbf{tecnologias} - Profissional do componente Matemática : 1 \textbf{aula} em todas as \textbf{matrizes}. - Área das Ciências da natureza e suas \textbf{tecnologias} - Profissional do componente Biologia ou Química : 1 \textbf{aula} ou 2 \textbf{aulas}, a \textbf{depender} da \textbf{matriz} em funcionamento na Unidade Escolar. 5.2.1 \textbf{Introdução} Que espaço é esse em que todos nós habitamos?
Há uma crescente urbanização do território brasileiro em relação às cidades do mundo quando \textbf{diz} respeito à divisão geográfica desigual e assimétrica dos espaços urbanos ( DAVIS, 2006 ).
As cidades de Santa Catarina, inclusive, estão em constante transformação desse espaço geográfico, ocupado de diferentes formas pela ação humana.
E, como tal, há múltiplas dimensões de seus habitantes no que \textbf{diz} respeito a questões de geração, classe, gênero, etnia, entre outras intersecções, para ocorrer sua produção.
Assim, cidade como forma de espaço é um conceito criado para abarcar este fato social urbano, pois não é possível pensar em sociedade separada do espaço que ocupa ; em outras palavras, ocupar o espaço é produto da ação humana e, como tal, tem historicidade. “ O espaço deve ser considerado como um conjunto indissociável de que participam, de um \textbf{lado}, \textbf{certo} arranjo de objetos geográficos, objetos naturais e objetos sociais, e, de outro, a vida que preenche e os anima, ou seja, a sociedade em movimento ” ( SANTOS, 1998, p. 26 ).
A fotografia, por sua vez, é integrada à discussão, como temática e como instrumento, pois, a partir dela, é possível trazer à tona memórias e narrativas dos \textbf{atores} sociais e, desta forma, ressignificar os sentidos construídos em torno das cidades.
Dobal ( 2017 ), ao ler Susan Sontag, conclui que a fotografia é, para além do uso instrumental, uma possibilidade de engajamento social.
Diante disso, a justificativa do tema “ Foto(cidades)grafias em movimento ” reside na complexidade da relação entre seus \textbf{atores}, ou grupos sociais, e espaços praticados das cidades a serem investigadas pelos(as) estudantes.
As cidades são espaços relacionais, de conexões, e exercem um papel social, ideológico e histórico no mundo.
Elas são capazes de \textbf{influenciar} outras, distantes de seu espectro.
Os habitantes das cidades exercem o poder de produzir, armazenar e distribuir produtos, serviços e informações.
Produzem estéticas plurais e espaços arquitetônicos, em que há diferenças de acesso e de uso entre os \textbf{atores} e grupos sociais das cidades.
Uma perspectiva dessa \textbf{trilha} é analisar as relações de poder \textbf{imbricadas} nas formas e funções arquitetônicas dos planejamentos urbanos como espaços praticados e \textbf{resultantes} das múltiplas condições de possibilidade histórico-socioculturais.
Nos espaços das cidades, como \textbf{categoria} dessa \textbf{trilha}, consideram -se os habitantes como narradores das suas experiências, na esteira de Michel de Certeau ( 1996 ).
Assim, há estéticas urbanas que se debruçam em duas redes de \textbf{análise} : os gestos e os relatos, ambas conectadas pelas formas de usos criativos desses objetos urbanos.
Alerta -se que os habitantes desfavorecidos nos espaços, como Michel de Certeau ( 1996, p. 198 ) defendeu, “ não só têm, no \textbf{quadro} da lei, o direito à ocupação dos lugares ; também têm direito à sua estética ”.
Compreendemos que a palavra estética — “ aisthesis ” — indica a \textbf{primordial} \textbf{capacidade} do ser humano de sentir a si próprio e ao mundo num todo integrado ( DUARTE JÚNIOR, 2011, p. 15 ).
Neste sentido, as cidades são pensadas pelas trajetórias singulares e experiências estéticas dos seus habitantes, e por isso passa pela leitura de sua dinâmica sociocultural.
Portanto, cidade é linguagem, signo e uma prática, e seria \textbf{explicitada} no processo de \textbf{ensino-aprendizagem} como “ a manifestação da linguagem do imaginário urbano.
São os gestos, as práticas, as artes de fazer e as narrativas do cotidiano ” ( DOSSE, 2004, p. 88 ).
Cabe -nos indicar as relações entre espaço e lugar.
Escreve Canton : “ A palavra ‘ espaço ’ é utilizada genericamente, enquanto lugar se refere a uma noção específica do espaço : trata -se de um espaço particular, familiar, responsável pela construção de nossas raízes e nossas referências no mundo ” ( 2009, p. 15 ).
Narrativas na e da cidade são lidas pelas lentes de Bakhtin ( 2000 ), ao considerar que a linguagem \textbf{sistematizada} nas narrativas não é neutra, apresentando caráter \textbf{dialógico} e de alteridade, sendo constituída por diversas \textbf{vozes} que perpassam as experiências dos sujeitos envolvidos na interação.
Neste sentido, a \textbf{trilha} trata dos espaços praticados das cidades interpretados por meio das narrativas, memórias e patrimônios, considerando suas disputas.
Primeiramente, as memórias dos \textbf{atores}, ou dos grupos sociais que compõem as narrativas das cidades, não correspondem a uma história desta, como contribuiu Pollack ( 1989 ) ao analisar que as memórias funcionam como fontes para a escrita da história.
Ele \textbf{explicitou} que nelas há esquecimentos, seletividades e silêncios.
Portanto, há, no tempo presente, o crescimento das narrativas memorialísticas acerca das experiências dos habitantes das cidades, e estas visam a construir sentidos para as práticas culturais, como um passado que persiste em ficar no presente.
Essa \textbf{compreensão} da memória se liga com a noção de patrimônio.
Este último, do ponto de \textbf{vista} do tempo histórico, é um sintoma de como lidar com as rupturas trazidas pelas grandes catástrofes do \textbf{século} XX.
Para reconhecer essas descontinuidades, a sociedade ocidental volta -se para estender o alcance da patrimonialização dos bens culturais de toda ordem.
Patrimônio é \textbf{tanto} integrante da vida cultural como instrumento das políticas públicas.
E, neste sentido, Vianna ( 2016 ) enunciou como possível patrimônio toda a cultura ocidental \textbf{moderna}, tendo no espectro um leque de coisas, bens de enorme valor para pessoas, comunidades, nações, ou toda a \textbf{humanidade}.
Os patrimônios se multiplicam nas cidades e fora dela.
Nesses usos e abusos de patrimonializar as cidades, cabe questionar : preservar o quê e quem?
Nas cidades do tempo presente há um processo de patrimonialização de todos os bens e práticas, que Hartog ( 2014 ) denomina de regime de historicidade presentista.
Nesse \textbf{debate} da patrimonialização acelerada dos bens culturais, que é atravessado por interesses políticos e \textbf{econômicos} da sociedade \textbf{contemporânea}, a noção de futuro tornou -se uma ameaça ( HARTOG, 2006 ).
Nesse ponto, o tempo passado é o receptáculo dos vestígios de identificações sociais que inventa “ lugares de memória ” ( NORA, 1993 ).
Nessa \textbf{trilha}, tem -se o alerta de Hartog ( 2006, p. 270 ) de o patrimônio se tornar “ um ramo principal da indústria do turismo ”, e, com isso, objeto de investimentos \textbf{econômicos} importantes.
Sua ‘ valorização ’ se insere, então, diretamente nos ritmos e temporalidades rápidas da economia de \textbf{mercado} de hoje, “ chocando -se e \textbf{aproximando} -se dela ”.
Nesse \textbf{debate}, há que se problematizar a patrimonialização do ambiente, questão que tem ocupado a Unesco, que usa dessa ação para preservar a diversidade cultural, a biodiversidade para o desenvolvimento durável, conceitos e objetivos que incidem em proteger o presente e preservar o futuro ( HARTOG, 2006 ).
Nessa chave de leitura, questiona -se : qual passado se congela no presente para lançar ao futuro?
A \textbf{trilha} apresenta as seguintes unidades curriculares : Cidades no tempo presente ; Cidades do bem viver coletivo ; Estética das arquiteturas ; Espaços democráticos da cidade ; Ateliê de fotografia.
Para \textbf{tanto}, operacionalizam -se as seguintes \textbf{categorias} como forma de ferramenta educativa : a ) conhecimento científico : desenvolver a \textbf{capacidade} de \textbf{compreensão} e reflexão com base em conhecimentos legitimados pela comunidade científica ; b ) cidades : como uma \textbf{categoria} e como temática, visa a investigar os espaços praticados ; c ) memórias : coletiva e individual, crivadas por condições de possibilidades do presente ao serem lembradas ; d ) narrativas : \textbf{categoria} que visa a perceber os modos de fazer e pensar dos habitantes das cidades e também como narrativas textuais, sonoras, visuais acerca das cidades ; e ) patrimônio : ligado à memória e ao território, que é uma ação no e sobre o presente. 5.2.2 Objetivo da \textbf{trilha} de \textbf{aprofundamento} Esta \textbf{trilha} de \textbf{aprofundamento} tem por objetivo dar luz ( foto ) à reflexão acerca das relações entre os espaços das cidades e seus \textbf{atores} sociais no território catarinense, analisando o movimento de patrimonialização, as narrativas, as memórias políticas e estéticas. 5.2.3 Unidades curriculares da \textbf{trilha} de \textbf{aprofundamento} - Unidade Curricular I : Cidades no tempo presente ( Eixos : Investigação Científica, Mediação e Intervenção sociocultural ) - Unidade Curricular II : Cidades do bem viver coletivo ( Eixos : Investigação Científica, Processos de Criação, Mediação e Intervenção sociocultural ) - Unidade Curricular III : Estética e arquitetura ( Eixos : Investigação Científica, Processos de Criação, Mediação e Intervenção sociocultural ) - Unidade Curricular IV : Espaços democráticos das cidades ( Eixos : Investigação Científica, Processos de Criação, Mediação e Intervenção sociocultural ) - Unidade Curricular V : Ateliê de fotografia ( Eixos : Investigação científica, Processos de criação, Mediação e Intervenção sociocultural, Empreendedorismo ) 5.2.4 Organização das unidades curriculares da \textbf{trilha} de \textbf{aprofundamento} Unidade Curricular I – Cidades no tempo presente Na unidade “ Cidades no tempo presente ”, são apresentados três objetos do conhecimento que visam a contribuir com o \textbf{aprofundamento} de discussões em torno de patrimônios históricos e culturais, memórias e lugares de vivência, a fim de confrontar narrativas históricas apresentadas e possibilitar ressignificações dos sentidos construídos em torno das cidades.
A unidade oportuniza aos(às) estudantes compreender as cidades no tempo presente por meio de memórias em confronto com narrativas históricas, e a \textbf{mobilizar} conceitos, conteúdos e conhecimentos com \textbf{vistas} à socialização, à atuação e à intervenção na comunidade.
A perspectiva \textbf{metodológica} dessa unidade considera os(as) estudantes como protagonistas das ações de ouvir, registrar e dar espaço a memórias e narrativas dos \textbf{atores} sociais.
As atividades propostas devem provocar investigações e experimentações próprias ( individuais ) e coletivas, colocando o estudante como protagonista de \textbf{reelaboração} de conceitos a partir do percurso realizado e de sua inserção em novas práticas em contexto escolar.
Ao reconhecer e valorizar o patrimônio histórico e cultural das cidades, \textbf{tanto} o patrimônio material quanto o imaterial, a intenção é que esse movimento possibilite diversas relações \textbf{críticas} em torno das cidades — composta de pessoas e pelas pessoas, inseridas em diversas práticas sociais, que, por sua vez, são \textbf{materializadas} nos modos de ser e de viver em diferentes tempos e espaços, \textbf{permeadas} por identidades e manifestações linguísticas dos diferentes grupos sociais.
Ao entrar em contato com esses patrimônios, que constituem as cidades no tempo presente, possibilita -se o contato com as memórias dos lugares de vivência, memórias essas que são constituídas por diferentes modos de vida dos lugares de vivência, seja no campo ou na cidade, e por um atravessamento de \textbf{vozes}, que também constituem as cidades.
Com relação às narrativas históricas, é válido destacar que a intenção, neste objeto de conhecimento, é registrar e \textbf{sistematizar} as memórias dos \textbf{atores} que integram o coletivo das cidades, de forma a \textbf{divulgá} -las na comunidade.
Neste sentido, tal exercício pode oportunizar a ressignificação dos sentidos que constituem os lugares de vivência, \textbf{tanto} para os(as) estudantes, quanto para as demais pessoas da comunidade e, ao mesmo tempo, dar \textbf{voz} e espaço de circulação aos \textbf{posicionamentos} silenciados e não legitimados. \textbf{Quadro} 44 – Cidades do tempo presente - Unidade Curricular : I — Cidades no tempo presente - Carga Horária : 40 h ou 60 h, a \textbf{depender} da \textbf{matriz} em funcionamento na Unidade Escolar - Habilidades da \textbf{trilha} associadas aos eixos estruturantes : - Selecionar e \textbf{sistematizar}, com base em estudos de campo, visando fundamentar reflexões e hipóteses sobre a organização, o funcionamento e/ou os efeitos de sentido de enunciados e discursos \textbf{materializados} nas diversas línguas e linguagens ( verbal, matemática, corporal, visual, sonora ou \textbf{digital} ) em torno das memórias e lugares de vivências, identificando os diversos pontos de \textbf{vista} e posicionando -se mediante argumentação, com o cuidado de citar as fontes dos recursos utilizados na pesquisa e buscando apresentar conclusões com o uso de diferentes mídias. - Selecionar e \textbf{mobilizar} intencionalmente conhecimentos e recursos das Ciências Humanas e Sociais Aplicadas para propor ações individuais e/ou coletivas de mediação e intervenção sobre \textbf{problemas} de natureza sociocultural relacionados aos sentidos construídos em torno dos lugares de memória nas cidades. - Propor e testar estratégias de mediação e intervenção sociocultural e ambiental, selecionando adequadamente elementos das diferentes linguagens, a fim de respeitar a diversidade dos \textbf{atores} sociais que constituem e se constituem nas cidades. - Registrar memórias dos lugares de vivência a fim de socializar, a partir de gêneros \textbf{discursivos}, as memórias investigadas e registradas, de \textbf{maneira} a dar \textbf{voz} e espaço de circulação aos \textbf{posicionamentos} muitas vezes silenciados ou não legitimados. - Habilidades da \textbf{trilha} associadas às áreas do conhecimento : - Identificar os conceitos que definem os patrimônios históricos e culturais, para compreender as diversas relações que estes estabelecem com espaços e lugares de vivência. - Reconhecer e refletir acerca dos patrimônios históricos e culturais \textbf{materializados} nas práticas sociais, nos lugares de vivência em múltiplos tempos e espaços, considerando as identidades dos diferentes grupos sociais. - Identificar as políticas públicas das cidades voltadas para os lugares de vivência. - Respeitar e valorizar as memórias das diversidades dos modos de vida nos lugares de vivência ( cidade e campo ), respeitando identidades. - Compreender as discussões e tensões em torno de memória, esquecimento e silêncio que \textbf{permeiam} os lugares de vivência em relação aos diferentes grupos sociais. - Elaborar narrativas históricas a partir das memórias dos diferentes grupos sociais acerca dos lugares de vivência, ressignificar os sentidos construídos em torno da cidade. - Investigar memórias dos lugares de vivência a fim de oportunizar reflexões em torno dos patrimônios materiais e imateriais que integram as cidades e o tempo presente. - Objetos de conhecimento : - Patrimônios históricos e culturais - Lugares de memória - Memórias dos lugares de vivência - Narrativas históricas Fonte : Elaboração dos \textbf{autores} ( 2020 ).
Unidade Curricular II – Cidades do bem viver coletivo Nesta unidade são propostos \textbf{debates} e reflexões acerca das narrativas das cidades em conexão com as disputas por lugares plurais e de ‘ bem viver ’ coletivo, considerando as diversidades ambientais.
As cidades ganham sentido de lugares de vivência, nos quais o ‘ bem viver ’ tem relação direta com a noção de natureza e ambiente como vital para a vida humana.
O sentido do ‘ bem viver ’ advém da raiz ancestral ligada à narrativa \textbf{intercultural} e das diversidades de modos de pensar no coletivo dos povos indígenas ( ACOSTA, 2016 ; SANTOS, 2010 ).
Esse referencial foi escolhido como forma de refletirmos a respeito das narrativas advindas dos movimentos ambientais em confronto com as políticas de gestão ambiental, bem como em relação às narrativas individualistas de globalização neoliberal, também vistos nas batalhas acerca do movimento de patrimonialização do ambiente.
Neste sentido, é \textbf{consenso} na produção bibliográfica acerca da noção de patrimônio ambiental, isso na esteira da nova forma de compreender o conceito de Natureza, ligado aos estudos culturais, “ uma vez que a fronteira entre Natureza e Cultura, constructo da modernidade europeia, tornou -se tênue e, para alguns, inexistente ” ( KARPINSKI, 2020, p. 315 ).
Ainda, conforme Hartog ( 2006 ), o movimento de patrimonialização ambiental “ designa a extensão provavelmente a mais massiva e a mais nova da noção, abre indubitavelmente sobre o futuro ou sobre novas interações entre presente e futuro ” ( p. 272 ).
Nesse ponto de \textbf{vista}, o futuro está ameaçado se não ocorrer a preservação do ambiente.
Assim, o sentido do objetivo desta Unidade Curricular é compreender as cidades como lugares de ‘ bem viver ’ coletivo, em relação \textbf{reflexiva} acerca dos movimentos de conservação e preservação do ambiente.
Questiona -se : o que preservar e por quê?
Para alcançar este objetivo foram recortados três objetos de conhecimento, que são : patrimônios naturais ou ambientais ; gestão ambiental ; movimentos ambientalistas.
A perspectiva \textbf{metodológica} dessa unidade enfoca o ateliê pedagógico, compreendido como \textbf{provocação} de atividades de experimentações individuais e coletivas.
Para \textbf{tanto}, é necessária a criação de um ambiente de ateliê em aprendizagem \textbf{horizontal} entre professor(es) e estudantes.
Para \textbf{tanto}, \textbf{sugere} -se que se realizem “ oficinas ” de fotografia dos lugares de vivência a respeito da temática dos patrimônios ambientais ( paisagens ).
Para isso, tende -se a tocar na relevância temática para os(as) estudantes. \textbf{Quadro} 45 – Cidades do bem viver coletivo - Unidade curricular : Cidades do bem viver coletivo - Carga Horária : 30 h ou 50 h, a \textbf{depender} da \textbf{matriz} em funcionamento na unidade escolar - Habilidades da \textbf{trilha} associadas aos eixos estruturantes : - Selecionar e \textbf{sistematizar}, com base em estudos de campo de patrimônio, visando fundamentar reflexões e hipóteses sobre os discursos e narrativas \textbf{materializados} nas diversas línguas e linguagens em torno dos lugares de memória dos povos indígenas e quilombolas, ribeirinhos, caiçaras, faxinalenses, caboclos(as), entre outros, em seus lugares de vivências, identificando os pontos de \textbf{vista} e posicionando -se mediante argumentação, citando as fontes dos recursos utilizados na pesquisa e buscando apresentar conclusões com o uso das mídias. - Selecionar e \textbf{mobilizar} intencionalmente conhecimentos e recursos das Ciências Humanas e Sociais Aplicadas e das Ciências da Natureza para propor ações individuais e/ou coletivas de mediação e intervenção sobre \textbf{problemas} relacionados aos sentidos construídos em torno do patrimônio natural ou ambiental. - Reconhecer produtos e/ou processos criativos da temática das diversidades dos patrimônios naturais ou ambientais por meio de fruição, vivências e reflexão \textbf{crítica} sobre obras ou \textbf{eventos} de diferentes práticas culturais, ampliando o repertório pessoal. - Propor reflexões éticas, estéticas, criativas para \textbf{problemas} reais relacionados ao tema dos patrimônios naturais ou ambientais situados em seus processos históricos e/ou culturais, em múltiplas escalas. - Habilidades da \textbf{trilha} associadas às áreas do conhecimento : - Identificar os patrimônios naturais ou ambientais de sua cidade como lugar de bem viver coletivo, considerando as diversidades culturais desses patrimônios. - Compreender a importância da água e da terra, em suas diferentes formas, para a vida e saúde humana e de seus grupos sociais, percebendo -as como patrimônios naturais ou ambientais da \textbf{humanidade}. - Identificar e refletir sobre os patrimônios naturais ou ambientais a partir da política de gestão ambiental acerca dos espaços das cidades onde habitam os povos ( indígenas — povos Guarani, kaingang e xokleng — quilombolas, ribeirinhos, caiçaras, faxinalenses, caboclos(as) e demais povos originários ), historicizando os lugares de vivência. - Identificar e refletir acerca dos planos de gestão ambiental de suas cidades e lugares de vivência. - Compreender os principais movimentos ambientalistas que atuam na proteção e conservação do ambiente, considerando as ações de conscientização a respeito do bem viver coletivo e de educação do patrimônio ambiental no território catarinense. - Identificar e refletir acerca das ações realizadas por movimentos ambientalistas em defesa do ambiente do território catarinense. - Objetos de conhecimento : - Patrimônios naturais ou ambientais - Gestão ambiental - Movimentos ambientalistas Fonte : Elaboração dos \textbf{autores} ( 2020 ).
Unidade Curricular III – Estéticas das arquiteturas Esta unidade curricular propõe a percepção estética das arquiteturas das cidades, considerando os diversos lugares que ocupamos, de que forma eles se apresentam no tempo e nos constituem como sujeitos.
A seguir, apresentam -se os seguintes objetos de conhecimento : Patrimônio Arquitetônico — considera -se a arquitetura como ofício ao longo do tempo histórico, como linguagem e como processo.
Ressignificar as particularidades estilísticas históricas, representacionais e \textbf{hegemônicas}, locais e regionais, é o primeiro passo para que se possa conduzir os(as) estudantes a mergulhar no universo estético e arquitetônico em que vivem e que vivenciam.
É \textbf{imperativo} que os(as) estudantes compreendam o universo sociocultural e a trajetória histórica do espaço em que está inserido o patrimônio arquitetônico para significar a estética da cidade no presente, no passado e, quiçá, no futuro, possibilitando a reflexão sobre o sentimento de pertencimento ao espaço ocupado.
Para \textbf{tanto}, \textbf{sugerem} -se como possibilidade \textbf{metodológica} : pesquisa sobre a história da cidade por meio de \textbf{entrevistas} dos seus habitantes, \textbf{debates} sobre as imagens da cidade projetadas pelo poder público e pelos(as) estudantes.
O segundo — Arquitetura Sensível — tem como princípio a \textbf{análise} da estética dos espaços por onde se passam e são constituídas as nossas \textbf{juventudes}.
As \textbf{análises} das experiências estéticas contam, recontam, atravessam -nos e nos constituem como \textbf{atores} sociais.
Ao aprofundar esta \textbf{análise} em relação às arquiteturas da cidade, os(as) estudantes precisam ser conduzidos a entender que vivem em espaços privados em meio à coisa pública.
Nesse espaço, nem tudo consegue estar planejadamente harmonizado.
Como metodologia, apontam -se : saídas a campo, para uma experiência estética intencional das cidades ; \textbf{exposição} oral dos sentimentos em relação às arquiteturas ; \textbf{aprofundamento} sobre o conceito de belo em sua relação com as arquiteturas.
Por fim, em Forma e Função da Arquitetura, pretende -se transcender a leitura estética em seus aspectos formais, relacionando -a com seus contextos de uso e significado, inclusive sobre as questões de acessibilidade a esses espaços.
Entende -se que os(as) estudantes precisam conhecer o que levou o ser humano a pensar nesta arquitetura que temos hoje, e sobre isso refletir, assim como a respeito da ocupação e a representação que fazem dos espaços públicos e privados.
É preciso que os(as) estudantes reconheçam nossas ruas como espaços de expressão, reflexão e transformação social, e compreendam as aproximações e distanciamentos entre os conceitos de Arte e Cultura como agentes de manutenção ou mudança das arquiteturas da cidade.
São metodologias : saídas a campo pela cidade, com olhares específicos para as estéticas da cidade em sua relação com as formas, os usos e os significados da cidade ; a identificação dos espaços de interlocuções artísticas na cidade ; o mapeamento de grafites e a aproximação de grafiteiros com a escola ; a pesquisa sobre as relações que fazem espaços e sujeitos receberem mais atenção do poder público em detrimento de outros. \textbf{Quadro} 46 – Estética das arquiteturas - Unidade Curricular : Estética das arquiteturas - Carga Horária : 30 h ou 50 h, a \textbf{depender} da \textbf{matriz} em funcionamento na unidade escolar - Habilidades da \textbf{trilha} associadas aos eixos estruturantes : - Selecionar e \textbf{sistematizar}, com base em estudos e/ou pesquisas em fontes confiáveis, informações sobre o patrimônio arquitetônico das cidades, visando fundamentar reflexões com o cuidado de citar as fontes dos recursos utilizados na pesquisa e buscando apresentar conclusões com o uso de diferentes mídias. - Reconhecer processos criativos por meio de fruição, vivências e reflexão \textbf{crítica} sobre as arquiteturas das cidades. - Propor soluções para \textbf{problemas} reais e/ou estratégias de mediação e intervenção sociocultural e ambiental, éticas, estéticas, criativas por meio da Arte ( forma e conteúdo ), opondo -se à estereotipia, ao lugar-comum e ao clichê. - Selecionar e \textbf{mobilizar} intencionalmente conhecimentos e recursos da Arte para propor ações individuais e/ou coletivas de mediação e intervenção sobre formas de interação e de atuação social, \textbf{artístico-cultural}, visando colaborar para o convívio democrático com a diversidade humana e ambiental. - Habilidades da \textbf{trilha} associadas às áreas do conhecimento : - Identificar as particularidades estilísticas históricas dos patrimônios arquitetônicos para compreender a participação dos \textbf{atores} ou grupos sociais na construção das identidades. - Analisar a arquitetura da cidade com um olhar estético, sensível e \textbf{crítico}, possibilitando a reflexão sobre o sentimento de pertencimento aos espaços ocupados e lugares de vivência. - Identificar e valorizar as manifestações artísticas e culturais que perpassam as arquiteturas das cidades. - Compreender a arquitetura inclusiva com respeito à acessibilidade aos lugares de vivência. - Identificar os fatores que \textbf{influenciam} as arquiteturas abandonadas e/ou destruídas e suas possíveis reutilizações. - Transcender os aspectos formais da leitura e fruição estética das arquiteturas, relacionando -os, criticamente, com os contextos de uso e significados. - Objetos de conhecimento : - Patrimônio arquitetônico - Arquitetura sensível - Forma e função da arquitetura Fonte : Elaboração dos \textbf{autores} ( 2020 ).
Unidade Curricular IV : Espaços democráticos da cidade Nesta unidade, apresentamos as cidades a partir do seu espaço e lugar, que pode ser físico, onde ocorrem as relações de trabalho, de permanência ou de pertencimento, mas também podem ser lugares afetivos, que envolvem memórias, nos quais a vida se desenvolve a partir de \textbf{inúmeras} conexões.
Cabe -nos indicar as relações entre espaço e lugar. “ A palavra ‘ espaço ’ é utilizada genericamente, enquanto lugar se refere a uma noção específica do espaço : trata -se de um espaço particular, familiar, responsável pela construção de nossas raízes e nossas referências no mundo ” ( CANTON, 2009, p. 15 ).
O objetivo geral desta unidade curricular é compreender a distribuição democrática dos espaços das cidades.
Os objetos de conhecimento propostos para essa unidade são : planejamento urbano e plano diretor ; espaços públicos, cartografia dos espaços coletivos e dimensões artísticas da cidade.
Estes objetos consideram as cidades a partir de elementos técnicos, mas também sensíveis, procurando \textbf{aproximar} pessoas a partir dos seus saberes e fazeres nesse espaço e lugar.
A perspectiva \textbf{metodológica} desta unidade pode se realizar por meio de pesquisas de campo e bibliográficas, envolvendo os(as) estudantes do ensino \textbf{médio} com a possibilidade de realizarem cartografias dos espaços públicos, procurando evidenciar os lugares democráticos da cidade e suas dimensões artísticas.
Cartografia é a \textbf{ciência} que estuda os mapas, mas também é uma metodologia de pesquisa, que \textbf{instiga} novas formas de perceber, pensar e fazer, que vão se construindo à medida que o/a estudante/cartógrafo vai assumindo sua posição de \textbf{pesquisador} e se dispõe a uma atitude investigativa, interessando -se mais pelos efeitos do que pelas causas.
Para cartografar, podem ser utilizados diferentes materiais, não privilegiando apenas os signos linguísticos, abrindo possibilidades para as visualidades, como fotografias, \textbf{vídeos}, pinturas e desenhos ; para as sonoridades, por meio de músicas, da paisagem sonora dos lugares de vivência, mas também a sonoridade que deixa transparecer a potência dos relatos obtidos por meio de \textbf{entrevistas} e rodas de conversa gravadas, quando as memórias coletadas trazem a força e a suavidade transmitida pela narrativa vocal, assim como o uso de outros materiais que possam ajudar no mapeamento das semióticas daquele campo. “ Cartografar é habitar um território existencial ” ( ALVAREZ e PASSOS, 2014, p. 131 ), o que indica a importância da imersão do estudante/cartógrafo no espaço da pesquisa, procurando compreendê -lo.
As atividades propostas devem provocar investigações e experimentações próprias, individuais e coletivas, nas quais o processo investigativo se fortaleça e obtenha relevância pessoal e coletiva. \textbf{Quadro} 47 – Espaços democráticos da cidade - Unidade Curricular : Espaços democráticos das cidades - Carga Horária : 30 h ou 50 h, a \textbf{depender} da \textbf{matriz} em funcionamento na unidade escolar - Habilidades da \textbf{trilha} associadas aos eixos estruturantes : - Realizar uma investigação científica por meio da \textbf{sistematização} da ressignificação das cidades no que se refere à sua identidade, histórias, participação e ocupação dos espaços públicos pelos cidadãos, promovendo a convivência e a troca de saberes. - Discutir a noção de ação performativa e material nos espaços públicos como planos de composição poética, por meio de processos criativos e relacionais que promovem maior intensidade na experiência coletiva. - Compreender a distribuição democrática dos espaços coletivos das cidades e lugares de vivência de acesso às dimensões artísticas e culturais, nos quais o empreendedorismo cultural e arquitetônico possam promover maior conectividade entre \textbf{atores} e grupos sociais, considerando a acessibilidade. - Habilidades da \textbf{trilha} associadas às áreas do conhecimento : - Reconhecer os espaços públicos das cidades e seus múltiplos usos. - Investigar o planejamento urbano e plano diretor das cidades com ênfase em espaços de vivência democráticos. - Produzir a cartografia dos espaços coletivos, refletindo acerca do acesso às dimensões artísticas das cidades. - Ocupar os espaços públicos a partir de possibilidades poéticas ( performativas e materiais ). - Reconhecer e utilizar variadas formas de representação do espaço : \textbf{cartográfica}, tratamentos gráficos, matemáticos, estatísticos, iconográficos. - Objetos de conhecimento : - Planejamento urbano e plano diretor - Espaços públicos - Cartografia dos espaços coletivos - Dimensões artísticas da cidade Fonte : Elaboração dos \textbf{autores} ( 2020 ).
Unidade Curricular V – Ateliê de fotografia Na unidade “ Ateliê de Fotografia ” são apresentados objetos de conhecimento que visam ao \textbf{aprofundamento} em torno da questão da fotografia e dos seus significados enquanto narrativas dos espaços de vivência.
Os objetos de conhecimento propostos consideram a pesquisa em fotografia, a composição e produção fotográfica e a curadoria em fotografia, elementos importantes na construção de conhecimentos necessários ao \textbf{aprofundamento} dos estudos pensados para essa \textbf{trilha}.
Dessa forma, ao estudar a pesquisa em fotografia, considera -se que os(as) estudantes conheçam a trajetória da fotografia em diferentes espaços e temporalidades, bem como seus usos e intenções, que fazem desse recurso tecnológico um grande veiculador de ideias.
Esse estudo permite “ pensar como as fotografias participam da construção da nossa imaginação — da realidade — do mundo \textbf{contemporâneo} ” ( OLIVEIRA JR., 2011, p. 257 ).
Ainda, permite aos(às) estudantes estabelecerem conexões entre as imagens fotográficas e a \textbf{capacidade} de evocarem memórias, sejam elas da vida individual, como da vida coletiva dos lugares em que estes vivem.
Nos estudos da composição e produção fotográfica, considera -se a técnica da organização e do enquadramento, do ponto de \textbf{vista}, das linhas, das cores, do equilíbrio destes fatores, do destaque do assunto principal, bem como da disposição dos assuntos secundários e dos recursos tecnológicos ( pré-produção, captura das imagens e posterior tratamento \textbf{digital} ) como elementos \textbf{indispensáveis} à construção das fotografias. \textbf{Espera} -se que os(as) estudantes conheçam esse processo, a fim de que compreendam que a técnica, aliada ao olhar do fotógrafo, é capaz de construir significados diferenciados, a partir de um recorte da mesma realidade, ou seja, a \textbf{maneira} como os lugares de vivências são retratados.
Permite a construção de uma imaginação espacial em que os mesmos acabam sendo colocados no patamar da riqueza ou da pobreza ( OLIVEIRA Jr., 2011 ).
Em curadoria em fotografia, os(as) estudantes terão a possibilidade de conhecer o processo que envolve a pesquisa, a seleção e a organização de mostras fotográficas.
Esse conhecimento \textbf{auxiliará} na elaboração da atividade proposta para essa Unidade Curricular ( uma \textbf{exposição} fotográfica ). \textbf{Quadro} 48 – Ateliê de fotografia - Unidade Curricular : Ateliê de Fotografia - Carga Horária : 30 h ou 50 h, a \textbf{depender} da \textbf{matriz} em funcionamento na unidade escolar - Habilidades da \textbf{trilha} associadas aos eixos estruturantes : - Investigar e analisar a organização, o funcionamento e/ou os efeitos de sentido de enunciados e discursos \textbf{materializados} nas fotografias, situando -os no contexto de um ou mais campos de atuação social dos sujeitos. - Investigar e analisar situações-problema envolvendo os temas sensíveis das imagens fotográficas que consideram processos históricos, sociais, \textbf{econômicos}, filosóficos, políticos e/ou culturais. - Reconhecer as fotografias e seus processos criativos por meio de fruição, vivências e reflexão \textbf{crítica} sobre obras ou \textbf{eventos} de diferentes práticas artísticas, culturais e/ou corporais, ampliando o repertório/domínio pessoal sobre o seu funcionamento. - Propor reflexões acerca das fotografias e seus temas, de forma ética, estética e criativa. - Identificar e explicar questões socioculturais e ambientais passíveis de mediação e intervenção por meio da fotografia. - Identificar e explicar situações em que ocorram conflitos, desequilíbrios e ameaças a grupos sociais, à diversidade de modos de vida, às diferentes identidades culturais e ao meio ambiente por meio das fotografias. - Propor fotografia como estratégia de mediação e intervenção para resolver \textbf{problemas} socioculturais e ambientais. - Habilidades da \textbf{trilha} associadas às áreas do conhecimento : - Pesquisar a trajetória da fotografia dos lugares de vivências, dos múltiplos tipos de fotografia ( publicitária, documental, jornalística, artística, didática ou científica ), e da questão ética do direito de imagem. - Conhecer e experimentar diferentes elementos da linguagem fotográfica, suas técnicas e seus processos compositivos. - Reconhecer os usos das fotografias como instrumentos que problematizam os lugares de memória e a patrimonialização histórica, cultural, ambiental, arquitetônica. - Reconhecer as narrativas de fotografias existentes dos \textbf{atores} e lugares de vivência, e produzir narrativas a partir de novas fotografias. - Produzir o gênero \textbf{discursivo} legenda a fim de contextualizar as fotografias. - Promover Mostra Fotográfica enfatizando a presença de valores sociais e humanos atualizáveis e permanentes no mundo fotográfico. - Objetos de conhecimento : - Pesquisa em fotografia - Composição e produção fotográfica - Curadoria em fotografia Fonte : Elaboração dos \textbf{autores} ( 2020 ). 5.2.5 Orientações \textbf{metodológicas} \textbf{Sugere} -se que o percurso \textbf{metodológico} desta \textbf{trilha} seja desenvolvido conforme \textbf{sequência} apresentada ; entretanto, entende -se que, de acordo com a realidade local, todo arranjo possa ser repensado para que a proposta melhor \textbf{atinja} seu objetivo.
O trabalho coletivo e colaborativo deve ser valorizado, \textbf{tanto} pelos docentes, no âmbito do planejamento com seus pares, quanto pelos(as) estudantes, inclusive em sua integração com a comunidade.
Sendo assim, parcerias são muito bem-vindas.
Além das parcerias, esta \textbf{trilha} requer estratégias de ensino e aprendizagem que priorizem estudos e pesquisas que identifiquem como, e se, os sujeitos das culturas escolares reconhecem as problemáticas que envolvem as cidades e seus lugares de vivência no território catarinense, promovendo : a investigação científica ; fóruns, \textbf{debates}, campanhas e \textbf{palestras} ; políticas públicas e projetos de intervenção urbana ; grupos de memória e de escuta dos moradores das cidades ; \textbf{entrevistas} com moradores de diferentes gerações ; rodas de conversa intergeracional nas escolas ou em lugares de vivência ; processos de criação e expressão artística ; mediação e intervenção sociocultural ; investigar as notícias \textbf{publicadas} na imprensa escrita local que tratem de questões relacionadas ao patrimônio cultural.
Usar nas \textbf{entrevistas} orais a metodologia da história oral com os “ profissionais do patrimônio ” — historiadores, arquitetos, urbanistas, geógrafos, arqueólogos, restauradores, educadores, juristas, dentre outros. \textbf{Entrevistas}, quando for o caso, com pessoas ligadas às \textbf{instâncias} consultivas e deliberativas, tipo Conselho Municipal de Política Cultural.
AVALIAÇÃO O processo \textbf{avaliativo}, que deve estar presente ao longo de todo o percurso da \textbf{trilha} integrada, considerando a multiplicidade de saberes a ela relacionados, assim como a diversidade de contextos e de sujeitos nos processos de ensinar e de aprender.
Na perspectiva da Proposta Curricular de Santa Catarina ( SANTA CATARINA, 2014 ) e do Currículo Base da Educação Infantil e do Ensino Fundamental do Território Catarinense ( SANTA CATARINA, 2019 ), a avaliação é concebida como investigação do processo de ensino e aprendizagem, possibilitando a tomada de decisão no desenvolvimento das atividades orientadoras de ensino.
Dessa forma, a avaliação aqui compreendida é fundamentada na legislação vigente em âmbito estadual e nacional ; deverá, portanto, estar explícita também no projeto político-pedagógico da escola, que deve apresentar a sua concepção de avaliação, a fim de contribuir para o planejamento e o replanejamento do trabalho docente. \textbf{Devido} ao caráter interdisciplinar da \textbf{trilha} de \textbf{aprofundamento}, os instrumentos de avaliação devem ser selecionados levando em consideração, sempre que possível, o diálogo entre as áreas de conhecimento envolvidas, a fim de contribuir no desenvolvimento e \textbf{aprofundamento} das \textbf{competências} e habilidades relacionadas às áreas e eixos dos itinerários formativos.
Dessa forma, os instrumentos, e respectivos critérios \textbf{avaliativos}, devem, em primeiro lugar, respeitar a diversidade dos(as) estudantes e das instituições escolares.
A seguir, devem ser apresentados e discutidos anteriormente à prática pedagógica, mantendo um diálogo com os(as) estudantes, de modo a considerar sua participação ativa, e, por fim, estar comprometidos com uma perspectiva formativa, aquela [... ] “ que ajuda o \textbf{aluno} a aprender a se desenvolver, ou melhor, que participa da regulação das aprendizagens e do desenvolvimento no sentido de um projeto educativo ” ( PERRENOUD, 1999, p. 28 ).

% matched lemmas: aluno, análise, aprofundamento, aproximar, artístico-cultural, atingir, ator, aula, autor, auxiliar, avaliativo, capacidade, cartográfico, categoria, certo, ciência, competência, compreensão, consenso, contemporâneo, crítico, debate, depender, devido, dialógico, digital, discursivo, divulgar, dizer, econômico, ensino-aprendizagem, entrevista, esperar, evento, explicitar, exposição, hegemônico, horizontal, humanidade, imbricar, imperativo, indispensável, influenciar, inglês, instigar, instância, intercultural, introdução, inúmero, juventude, lado, maneira, materializar, matriz, mercado, metodológico, mobilizar, moderno, médio, palestra, permear, pesquisador, posicionamento, primordial, problema, provocação, publicar, quadro, reelaboração, reflexivo, resultante, sequência, sistematizar, sistematização, sugerir, século, tanto, tecnologia, trilha, vista, voz, vídeo
\end{textsample}
